\begin{abstract}
Nocturnal stable boundary layers (SBLs) exhibit complex interactions between three dominant scales: large-scale mesoscale pressure forcing, horizontally propagating internal gravity waves (IGWs), and highly localized turbulent shear bursts. Traditional diagnostic frameworks conflate these regimes, producing high variance in inversion-layer flux estimates and inconsistent model-observation comparisons. We present a \emph{metric-consistent Riemannian pseudospectral} framework that projects sparse tower observations onto a non-uniform computational grid designed explicitly to maximize sensitivity near the surface (\SI{\DzMinMeters}{\meter} effective node resolution at $z = \SI{\TowerMinMeters}{\meter}$) while maintaining numerical stability at upper boundaries ($\kappa(\mathbf{M}) = \KappaMassBaseline$). Using Chebyshev polynomials as the spectral basis and solving the rank-deficient least-squares problem via thin singular value decomposition (SVD), we extract three physically meaningful diagnostic features: effective modal dimension $D_{\mathrm{eff}}$, wave energy fraction $F_W$, and spectral curvature $\chi_N$.

Applying this framework to the CASES-99 campaign stable window (22--31 October 1999), we identify \emph{three statistically separable regimes} with a silhouette clustering score of $\overline{S} = \SilhouettePeak$ for $K=3$. Recent campaign diagnostics yield regime-mean effective dimensions of $D_{\mathrm{eff}}=\RegimeOneMean$ (continuous), $\RegimeTwoMean$ (wave-dominated), and $\RegimeThreeMean$ (intermittent), with windowed separability values of $\overline{S}_{\mathrm{early}}=\WindowSEarly$, $\overline{S}_{\mathrm{trans}}=\WindowSTransitional$, and $\overline{S}_{\mathrm{IOP}}=\WindowSIOP$. Because these means are partially collared across windows, regime attribution is drawn from the joint diagnostic manifold $[D_{\mathrm{eff}}, F_W, \chi_N, \mathrm{Ri}_g]$ and temporal persistence rather than fixed $D_{\mathrm{eff}}$ thresholds alone. Crucially, the regime-specific rank compression ($\Delta_{\mathrm{rank}} = \RankCompression$) confirms that atmospheric dynamics naturally occupy lower-dimensional manifolds and are not mathematical artifacts of basis truncation.

We extend the theoretical framework by establishing an explicit physical link to the Ozmidov buoyancy scale $L_O$, resolving the grid-topology mismatch via a generalized metric-consistent Virtual Tower interpolation operator ($\mathcal{P}_{\mathrm{VT}}$) for direct NWP model validation, and formulating a $2/3$-rule spectral anti-aliasing filter for forward-predictive implementations. A synthetic gravity-wave reflection test confirms that the spectral partition functions suppress upper-boundary reflection by approximately $3.3\times 10^4$, validating the sponge-layer boundary conditions.
\end{abstract}